% GENERATED FILE. Edit canonical lessons, then run npm run generate:books.
% canonical-source-hash: fnv1a64:719e58a2a5d3119e
% canonical-lessons: JA-C95-kanou, JA-C95-tokui, JA-C95-nigate, JA-C95-mochimono, JA-C95-okanemochi

\chapter{Possible, Good at, Poor at, Belongings, Rich person}
\label{ch:ja-possible-good-at-poor-at-belongings-rich-person}
\hlchaptermodality{95}

\begin{chapteropening}
\textbf{By the end of this chapter:} \emph{I can say possible, good at, poor at, belongings and a rich person.}

Complete the last lesson of chapter 95: I can say possible, good at, poor at, belongings and a rich person.
\end{chapteropening}

\section[\texorpdfstring{\ja{かのう} (kanou)}{かのう (kanou)}]{\ja{かのう} (kanou) — possible}

\label{lesson:JA-C95-kanou}



\begin{quote}
Before the new one: say the Japanese for something to look forward to, then the Japanese for love.
\end{quote}

\subsection*{You'll want to know: \ja{かのう}}
\textbf{\ja{かのう}} — \emph{kanou} — \textquotedblleft{}possible\textquotedblright{}.

\begin{grammarlens}[title={What you've built}]
One more word for this chapter.
\end{grammarlens}

\subsection*{Your turn}
\begin{itemize}
\raggedright
  \item \emph{Say it:} \emph{kanou}
  \item \emph{Say it:} \emph{kanou}, once more
  \item \emph{Say it:} say \emph{kanou}
  \item \emph{Recall:} say the Japanese for something to look forward to, then the Japanese for love, then say \emph{kanou} again
\end{itemize}

\begin{tcolorbox}[breakable,colback=teal!4,colframe=teal!35!black,title={Before you move on}]
What does \ja{かのう} mean? (\textquotedblleft{}Possible\textquotedblright{}.) Say it once more.
\end{tcolorbox}
\section[\texorpdfstring{\ja{とくい} (tokui)}{とくい (tokui)}]{\ja{とくい} (tokui) — good at}

\label{lesson:JA-C95-tokui}



\begin{quote}
Before the new one: say the Japanese for love, then the Japanese for possible.
\end{quote}

\subsection*{You'll want to know: \ja{とくい}}
\textbf{\ja{とくい}} — \emph{tokui} — \textquotedblleft{}good at\textquotedblright{}.

\begin{grammarlens}[title={What you've built}]
One more word for this chapter.
\end{grammarlens}

\subsection*{Your turn}
\begin{itemize}
\raggedright
  \item \emph{Say it:} \emph{tokui}
  \item \emph{Say it:} \emph{tokui}, once more
  \item \emph{Say it:} say \emph{tokui}
  \item \emph{Recall:} say the Japanese for love, then the Japanese for possible, then say \emph{tokui} again
\end{itemize}

\begin{tcolorbox}[breakable,colback=teal!4,colframe=teal!35!black,title={Before you move on}]
What does \ja{とくい} mean? (\textquotedblleft{}Good at\textquotedblright{}.) Say it once more.
\end{tcolorbox}
\section[\texorpdfstring{\ja{にがて} (nigate)}{にがて (nigate)}]{\ja{にがて} (nigate) — poor at}

\label{lesson:JA-C95-nigate}



\begin{quote}
Before the new one: say the Japanese for possible, then the Japanese for good at.
\end{quote}

\subsection*{You'll want to know: \ja{にがて}}
\textbf{\ja{にがて}} — \emph{nigate} — \textquotedblleft{}poor at\textquotedblright{}.

\begin{grammarlens}[title={What you've built}]
One more word for this chapter.
\end{grammarlens}

\subsection*{Your turn}
\begin{itemize}
\raggedright
  \item \emph{Say it:} \emph{nigate}
  \item \emph{Say it:} \emph{nigate}, once more
  \item \emph{Say it:} say \emph{nigate}
  \item \emph{Recall:} say the Japanese for possible, then the Japanese for good at, then say \emph{nigate} again
\end{itemize}

\begin{tcolorbox}[breakable,colback=teal!4,colframe=teal!35!black,title={Before you move on}]
What does \ja{にがて} mean? (\textquotedblleft{}Poor at\textquotedblright{}.) Say it once more.
\end{tcolorbox}
\section[\texorpdfstring{\ja{もちもの} (mochimono)}{もちもの (mochimono)}]{\ja{もちもの} (mochimono) — belongings}

\label{lesson:JA-C95-mochimono}



\begin{quote}
Before the new one: say the Japanese for good at, then the Japanese for poor at.

\emph{Read it:} the three greetings on sight

\begin{itemize}
\raggedright
  \item \textbf{\ja{こんにちは}}
  \item \textbf{\ja{ありがとう}}
  \item \textbf{\ja{さようなら}}
\end{itemize}
\end{quote}

\subsection*{You'll want to know: \ja{もちもの}}
\textbf{\ja{もちもの}} — \emph{mochimono} — \textquotedblleft{}belongings\textquotedblright{}.

\begin{grammarlens}[title={What you've built}]
One more word for this chapter.
\end{grammarlens}

\subsection*{Your turn}
\begin{itemize}
\raggedright
  \item \emph{Say it:} \emph{mochimono}
  \item \emph{Say it:} \emph{mochimono}, once more
  \item \emph{Say it:} say \emph{mochimono}
  \item \emph{Recall:} say the Japanese for good at, then the Japanese for poor at, then say \emph{mochimono} again
\end{itemize}

\begin{tcolorbox}[breakable,colback=teal!4,colframe=teal!35!black,title={Before you move on}]
What does \ja{もちもの} mean? (\textquotedblleft{}Belongings\textquotedblright{}.) Say it once more.
\end{tcolorbox}
\section[\texorpdfstring{\ja{おかねもち} (okanemochi)}{おかねもち (okanemochi)}]{\ja{おかねもち} (okanemochi) — a rich person}

\label{lesson:JA-C95-okanemochi}



\begin{quote}
Before the new one: say the Japanese for poor at, then the Japanese for belongings.
\end{quote}

\subsection*{You'll want to know: \ja{おかねもち}}
\textbf{\ja{おかねもち}} — \emph{okanemochi} — \textquotedblleft{}a rich person\textquotedblright{}.

\begin{grammarlens}[title={What you've built}]
One more word for this chapter.
\end{grammarlens}

\subsection*{Your turn}
\begin{itemize}
\raggedright
  \item \emph{Say it:} \emph{okanemochi}
  \item \emph{Say it:} \emph{okanemochi}, once more
  \item \emph{Say it:} say \emph{okanemochi}
  \item \emph{Recall:} say the Japanese for poor at, then the Japanese for belongings, then say \emph{okanemochi} again
\end{itemize}

\begin{tcolorbox}[breakable,colback=teal!4,colframe=teal!35!black,title={Before you move on}]
What does \ja{おかねもち} mean? (\textquotedblleft{}A rich person\textquotedblright{}.) Say it once more.
\end{tcolorbox}
